\documentclass[10pt,a4paper]{amsart}
\usepackage[margin=19mm]{geometry}
\usepackage{amsmath,amssymb,mathtools}
\usepackage[scr=rsfs,cal=euler]{mathalfa}
\usepackage{xfrac}
\usepackage[unicode,hidelinks,bookmarks=false]{hyperref}
\newcommand{\ie}{i.e.\ }
\newcommand{\cf}{cf.\ }
\newcommand{\resp}{respectively\ }
\newcommand{\Subsection}{\subsection}
\providecommand{\nobreakdash}{\nobreak\hbox{-}}
\setlength{\emergencystretch}{2em}
\multlinegap0pt
\usepackage{bm}
\usepackage{bbm}
\usepackage{dsfont}
\usepackage{xspace}
\usepackage{cancel}
\usepackage{mathtools}
\usepackage{amsthm}
\makeatletter
\@ifundefined{theorem}{\newtheorem{theorem}{Theorem}}{}
\makeatother
\title[Existence and nonexistence of spherical $5$-designs of minim]{Existence and nonexistence of spherical $5$-designs of minimal type}
\author{Sho Suda, Zili Xu, Wei-Hsuan Yu}
\date{Source: arXiv:2508.18685v1 (CC BY 4.0)}
\begin{document}
\maketitle

\noindent\emph{Source and licence.} Existence and nonexistence of spherical $5$-designs of minimal type. Version arXiv:2508.18685v1 (CC BY 4.0), distributed under the Creative Commons Attribution 4.0 International licence, as recorded in the arXiv licence metadata for this exact version and shown on the abstract page https://arxiv.org/abs/2508.18685v1. Full rights notice: licenses/arxiv-2508.18685-CC-BY-4.0.txt.\par\smallskip

\noindent\textit{Editorial scope.} Counted unit: the Theorem whose source label is \texttt{thm:di} in the article (source0.tex lines 449--458, source label \texttt{thm:di}), delivered together with the article's own complete original proof of it (source0.tex lines 459--505, one \texttt{proof} environment of 47 lines). No mathematics is written, repaired or re-derived: statement and proof are the article's own text line for line. Required context retained uncounted: sperical5design-2 (lines 680--683); sperical5design-4 (lines 680--683); sum (lines 691--693). Every definition, display and label of those lines is retained unchanged. Omitted from this package and not redistributed: 1--448, 506--679, 684--690, 694--1337. Nothing inside a retained excerpt is omitted or altered: every retained body is delivered exactly as the article's lines are written, so no omission receipt of any kind accompanies this package. Citations: the retained excerpts carry no citation command at all, so no bibliography entry of the article is transcribed and no reference is redistributed. Reference adaptation: none was needed and none was made; every reference inside the delivered unit resolves inside it, so no unresolved reference or citation is produced. Printed slips kept literal: no printed word, symbol, hypothesis or number of the counted statement or of its proof was altered. Layout and class: the article's own class is \texttt{article} with packages including amsmath, amsfonts, amsthm, amssymb, authblk, color, fullpage, url, graphicx, enumerate, dsfont, booktabs, mathrsfs, verbatim; the wrapper is the lane's pinned \texttt{amsart} editorial wrapper at 10pt on A4, which restates the theorem-like environments the retained text needs and adds the article's own macro definitions that the retained text needs (none), with extra wrapper packages (bm, bbm, dsfont, xspace, cancel, mathtools). The delivered numbering is the unit's own. Assets: no retained span contains a figure environment, a graphics-inclusion command, a PDF-page inclusion, a table or a TikZ picture; no image file is copied into this package, the package declares no asset dependency and no image file is redistributed with it. Licence: version arXiv:2508.18685v1 of this article is distributed under the Creative Commons Attribution 4.0 International licence, as recorded in the arXiv licence metadata for this exact version and shown on the abstract page https://arxiv.org/abs/2508.18685v1. A full rights notice is delivered at licenses/arxiv-2508.18685-CC-BY-4.0.txt. Characters: every retained excerpt is pure ASCII, so the delivered engine, XeLaTeX with its default Latin Modern text font, reports no missing character.
\begin{align}
\sum_{x\in D}\langle y,x \rangle^2&=\frac{1}{d}\cdot |D|\cdot \langle y,y \rangle,\label{sperical5design-2} \\
\sum_{x\in D}\langle y,x \rangle^4&=\frac{3}{d(d+2)}\cdot |D|\cdot\langle y,y \rangle^2 \label{sperical5design-4}.
\end{align}

\begin{equation}\label{sum}
n_0(\alpha)+n_1(\alpha)=|D|=d(d+1),
\end{equation}

\begin{theorem}\label{thm:di}
    Let $X_i$ be a spherical $t_i$-design in $S^{d-1}$ for $i\in\{1,\ldots,n\}$.  
    Let $s_{i,j}=|A(X_i,X_j)|$. 
    Assume that one of the following holds for $i,j\in\{1,\ldots,n\}$:
    \begin{enumerate}
        \item[{\rm (i)}] $s_{i,j}-1\leq t_j$,
        \item[{\rm (ii)}] $s_{i,j}-2= t_j$ and $X_i=-X_j$. Write $A'(X_i,X_j)=A(X_i,X_j)\setminus\{-1\}$. 
    \end{enumerate}
    Then the sets $\cup_{i=1}^n X_i$ is distance invariant.
\end{theorem}

\begin{proof}
    For $x\in X_i$ and $\alpha\in A(X_i,X_j)$, define $p_{\alpha}^j(x)=|\{z\in X_j \mid \langle x,z\rangle=\alpha \}|$. 
    We reprove equalities \eqref{sum}, \eqref{sperical5design-2}, \eqref{sperical5design-4} and prove them for exponent $\ell$ odd. 
    Let $H_\ell$ be a characteristic matrix of $X_j$. 
    Calculate $(\sum_{\ell=0}^\lambda f_{\lambda,\ell}\phi_{\ell}(x)H_{\ell}^\top)H_0$ in two ways. 

    On the one hand, 
    \begin{align}
        (\sum_{\ell=0}^\lambda f_{\lambda,\ell}\phi_{\ell}(x)H_{\ell}^\top)H_0&=\sum_{\ell=0}^\lambda f_{\lambda,\ell}\phi_{\ell}(x)H_{\ell}^\top H_0\nonumber\\
        &=f_{\lambda,0}\phi_{0}(x)H_{0}^\top H_0\nonumber\\
        &=|X_j|f_{\lambda,0},\label{eq:di1} 
    \end{align}
    and on the other hand, 
    \begin{align}
        (\sum_{\ell=0}^\lambda f_{\lambda,\ell}\phi_{\ell}(x)H_{\ell}^\top)H_0&=\sum_{z\in X_j}\sum_{\ell=0}^\lambda f_{\lambda,\ell}\phi_{\ell}(x)\phi_{\ell}(z)^\top \nonumber\\
        &=\sum_{z\in X_j}\sum_{\ell=0}^\lambda f_{\lambda,\ell}Q_{\ell}(\langle x,z\rangle)\nonumber\\
        &=\sum_{z\in X_j}\langle x,z\rangle^\lambda\nonumber\\
        &=\sum_{\alpha\in A(X_i,X_j)}\alpha^\lambda p_{\alpha}^j(x)+\delta_{X_i,X_j} \label{eq:di2}\\
        &=\sum_{\alpha\in A'(X_i,X_j)}\alpha^\lambda p_{\alpha}^j(x)+\delta_{X_i,X_j}+\delta_{X_i,-X_j}(-1)^\lambda. \label{eq:di3}
    \end{align}
    We consider the cases (i) and (ii) separately. 
    
    (i) We obtain from \eqref{eq:di1} and \eqref{eq:di2}:   
    \begin{align}\label{eq:di4}
        \sum_{\alpha\in A(X_i,X_j)}\alpha^\lambda p_{\alpha}^j(x)=|X_j|f_{\lambda,0}-\delta_{X_i,X_j}.
    \end{align}
    For $\lambda\leq s_{i,j}-1$, \eqref{eq:di4} yields  a system of $s_{i,j}$ linear equations whose unknowns are
    \[
    \{p_{\alpha}^j(x) \mid \alpha\in A(X_i,X_j)\}. 
    \]
    Its coefficients matrix is 
 the Vandermonde matrix $(\alpha^\lambda)$ where $\alpha\in A(X_i,X_j)$ and $\lambda\in\{0,\ldots,s_{i,j}-1\}$. 
 Therefore $p_{\alpha}^j(x)$ does not depend on the choice of $x\in X_i$, and is uniquely determined by $|X_j|$ and $A(X_i,X_j)$. 
 
    (ii) We obtain from \eqref{eq:di1} and \eqref{eq:di3}:   
    \begin{align}\label{eq:di5}
        \sum_{\alpha\in A'(X_i,X_j)}\alpha^\lambda p_{\alpha}^j(x)=|X_j|f_{\lambda,0}-\delta_{X_i,X_j}-(-1)^\lambda.
    \end{align}
    For $\lambda\leq s_{i,j}-2$, \eqref{eq:di5} yields a system of $s_{i,j}-1$ linear equations whose unknowns are
    \[
    \{p_{\alpha}^j(x) \mid \alpha\in A'(X_i,X_j)\}. 
    \]
    Its coefficients matrix is 
 the Vandermonde matrix $(\alpha^\lambda)$ where $\alpha\in A(X_i,X_j)$ and $\lambda\in\{0,\ldots,s_{i,j}-2\}$. 
 Therefore $p_{\alpha}^j$ does not depend on the choice of $x\in X_i$, and is uniquely determined by $|X_j|$ and $A'(X_i,X_j)$.
 Thus, a collection of $X_1,\ldots,X_n$ is distance invariant. 
\end{proof}

\end{document}
